\documentclass[10pt,a4paper]{amsart}
\usepackage[margin=19mm]{geometry}
\usepackage{amsmath,amssymb,mathtools}
\usepackage[scr=rsfs,cal=euler]{mathalfa}
\usepackage{xfrac}
\usepackage[unicode,hidelinks,bookmarks=false]{hyperref}
\newcommand{\ie}{i.e.\ }
\newcommand{\cf}{cf.\ }
\newcommand{\resp}{respectively\ }
\newcommand{\Subsection}{\subsection}
\providecommand{\nobreakdash}{\nobreak\hbox{-}}
\setlength{\emergencystretch}{2em}
\multlinegap0pt
\usepackage{bm}
\usepackage{bbm}
\usepackage{dsfont}
\usepackage{xspace}
\usepackage{cancel}
\usepackage{mathtools}
\usepackage{cleveref}
\usepackage{amsthm}
\makeatletter
\@ifundefined{theorem}{\newtheorem{theorem}{Theorem}}{}
\@ifundefined{proposition}{\newtheorem{proposition}[theorem]{Proposition}}{}
\makeatother
\newcommand{\mf}{\mathfrak}
\title[A curiously slowly mixing Markov chain]{A curiously slowly mixing Markov chain}
\author{Persi Diaconis, Andrew Lin, Arun Ram}
\date{Source: arXiv:2511.01245v2 (CC BY 4.0)}
\begin{document}
\maketitle

\noindent\emph{Source and licence.} A curiously slowly mixing Markov chain. Version arXiv:2511.01245v2 (CC BY 4.0), distributed under the Creative Commons Attribution 4.0 International licence, as recorded in the arXiv licence metadata for this exact version and shown on the abstract page https://arxiv.org/abs/2511.01245v2. Full rights notice: licenses/arxiv-2511.01245-CC-BY-4.0.txt.\par\smallskip

\noindent\textit{Editorial scope.} Counted unit: the Proposition whose source label is \texttt{closedformk} in the article (source0.tex lines 391--396, source label \texttt{closedformk}), delivered together with the article's own complete original proof of it (source0.tex lines 397--421, one \texttt{proof} environment of 25 lines). No mathematics is written, repaired or re-derived: statement and proof are the article's own text line for line. Required context retained uncounted: burnsidegenerictransition (lines 99--101). Every definition, display and label of those lines is retained unchanged. Omitted from this package and not redistributed: 1--98, 102--390, 422--903. Nothing inside a retained excerpt is omitted or altered: every retained body is delivered exactly as the article's lines are written, so no omission receipt of any kind accompanies this package. Citations: the retained excerpts carry no citation command at all, so no bibliography entry of the article is transcribed and no reference is redistributed. Reference adaptation: none was needed and none was made; every reference inside the delivered unit resolves inside it, so no unresolved reference or citation is produced. Printed slips kept literal: no printed word, symbol, hypothesis or number of the counted statement or of its proof was altered. Layout and class: the article's own class is \texttt{article} with packages including geometry, xcolor, tikz, tgpagella, graphicx, amsmath, amssymb, amsthm, mathtools, bbm, hyperref, cleveref, tocbibind; the wrapper is the lane's pinned \texttt{amsart} editorial wrapper at 10pt on A4, which restates the theorem-like environments the retained text needs and adds the article's own macro definitions that the retained text needs (\texttt{\textbackslash mf}), with extra wrapper packages (bm, bbm, dsfont, xspace, cancel, mathtools, cleveref). The delivered numbering is the unit's own. Assets: no retained span contains a figure environment, a graphics-inclusion command, a PDF-page inclusion, a table or a TikZ picture; no image file is copied into this package, the package declares no asset dependency and no image file is redistributed with it. Licence: version arXiv:2511.01245v2 of this article is distributed under the Creative Commons Attribution 4.0 International licence, as recorded in the arXiv licence metadata for this exact version and shown on the abstract page https://arxiv.org/abs/2511.01245v2. A full rights notice is delivered at licenses/arxiv-2511.01245-CC-BY-4.0.txt. Characters: every retained excerpt is pure ASCII, so the delivered engine, XeLaTeX with its default Latin Modern text font, reports no missing character.
\begin{equation}\label{burnsidegenerictransition}
    K(x, y) = \frac{1}{|G_x|} \sum_{s \in G_x \cap G_y} \frac{1}{|\mf{X}_s|}.
\end{equation}

\begin{proposition}\label{closedformk}
Let $K$ denote the transition matrix for the binary Burnside process. Fix $x, y \in C_2^n$, and for $a, b \in \{0, 1\}$, let $n_{ab}$ be the number of coordinates $i$ where $x_i = a$ and $y_i = b$ (so $n_{00} + n_{01} + n_{10} + n_{11} = n$). Then 
\[
    K(x, y) = \frac{\binom{2n_{00}}{n_{00}} \binom{2n_{01}}{n_{01}} \binom{2n_{10}}{n_{10}} \binom{2n_{11}}{n_{11}}}{4^n\binom{n_{00} + n_{01}}{n_{00}}\binom{n_{10} + n_{11}}{n_{10}}}.
\]
\end{proposition}

\begin{proof}
The permutations that fix $x$ can be described by $G_x = S_{n_{00} + n_{01}} \times S_{n_{10} + n_{11}}$, in which we permute the indices where $x_i = 0$ and also the indices where $x_i = 1$. Similarly, $G_x \cap G_y = S_{n_{00}} \times S_{n_{01}} \times S_{n_{10}} \times S_{n_{11}}$ is the set of all permutations whose cycles are each contained entirely within each type of coordinate. If $\sigma = \sigma_{00} \times \sigma_{01} \times \sigma_{10} \times \sigma_{11} \in G_x \cap G_y$, then
\[
    \frac{1}{|\mf{X}_g|} = \left(\frac{1}{2}\right)^{C(\sigma_{00}) + C(\sigma_{01}) + C(\sigma_{10}) + C(\sigma_{11})},
\]
where $C(\tau)$ denotes the number of cycles in the permutation $\tau$, since a binary $n$-tuple is fixed by $\sigma$ if and only if it is constant (either all $0$ or all $1$) on each cycle. Therefore, by plugging into \cref{burnsidegenerictransition},
\begin{align*}
    K(x, y) &= \frac{1}{(n_{00} + n_{01})! (n_{10} + n_{11})!} \sum_{\sigma_{00}, \sigma_{01}, \sigma_{10}, \sigma_{11}}\left(\frac{1}{2}\right)^{C(\sigma_{00}) + C(\sigma_{01}) + C(\sigma_{10}) + C(\sigma_{11})} \\
    &= \frac{1}{(n_{00} + n_{01})! (n_{10} + n_{11})!} \sum_{\sigma_{00}} \left(\frac{1}{2}\right)^{C(\sigma_{00})} \sum_{\sigma_{01}} \left(\frac{1}{2}\right)^{C(\sigma_{01})} \sum_{\sigma_{10}} \left(\frac{1}{2}\right)^{C(\sigma_{10})} \sum_{\sigma_{11}} \left(\frac{1}{2}\right)^{C(\sigma_{11})}.
\end{align*}
Now because the generating function for permutation cycle count is given by
\[
    C_n(x) = \sum_{\sigma \in S_n} x^{C(\sigma)} = x(x+1)\cdots(x+n-1)
\]
(for example by induction, since there are $(n-1)$ ways to insert the number $n$ into an existing cycle and $1$ way to add a new cycle with just $n$), we have that 
\[
    \sum_{\sigma_{00}} \left(\frac{1}{2}\right)^{C(\sigma_{00})} = \frac{1}{2} \left(\frac{1}{2} + 1\right) \cdots \left(\frac{1}{2} + n_{00} - 1\right) = \frac{(2n_{00}-1)!!}{2^{n_{00}}} = \frac{(2n_{00})!}{4^{n_{00}} n_{00}!}
\]  
and similarly for the other terms. Thus
\begin{align*}
K(x, y) &= \frac{1}{(n_{00} + n_{01})! (n_{10} + n_{11})!}  \frac{(2n_{00})!}{4^{n_{00}} n_{00}!} \frac{(2n_{01})!}{4^{n_{01}} n_{10}!} \frac{(2n_{10})!}{4^{n_{10}} n_{10}!} \frac{(2n_{11})!}{4^{n_{11}} n_{11}!} \\
&= \frac{1}{4^n} \frac{n_{00}! n_{01}! n_{10}! n_{11}!}{(n_{00} + n_{01})! (n_{10} + n_{11})!} \frac{(2n_{00})!}{n_{00}!^2} \frac{(2n_{01})!}{n_{01}!^2} \frac{(2n_{10})!}{n_{10}!^2} \frac{(2n_{11})!}{n_{11}!^2},
\end{align*}
which rearranges to the desired result.
\end{proof}

\end{document}
